\documentclass[11pt]{article}
\usepackage{fontspec,amsmath,amssymb,amsthm,hyperref}
\usepackage[margin=1in]{geometry}
\setlength{\emergencystretch}{3em}
\title{SGLD convergence for the actual control-variate chain}
\author{LangevinLean mathematical exposition}
\date{}
\begin{document}\maketitle
\noindent Source module: \texttt{AverageLMC/SGLDTheorem.lean}. Compile this proof document with XeLaTeX or LuaLaTeX. The accompanying HTML pairs the mathematical argument with the exact Lean statements and records the scope of the formalization.


The stochastic inputs are the component potentials and the independent Gaussian/index transition, with the anchor equal to the true minimizer. No transport recurrence or invariant law is supplied as a premise.

\[0\le h\le\frac{\bar m}{3nM_{\rm rms}^2}\quad\Longrightarrow\quad W_2^2(\nu_k,\pi)\le(1+n\bar m h)^{-k}W_2^2(\nu_0,\pi)+\left(M_{\rm av}(f/n)+2\bar m+\frac{2M_{\rm rms}^2}{\bar m}\right)\frac{hp}{\bar m}.\]
\section{Derive the unrestricted one-step claim}


Combine the actual average-smoothness LMC inequality with variance addition. The exact variance identity subtracts one full-gradient energy from estimator energy. The leftover coefficient is at most h squared times (hM-2)/(1+mh), because the discarded term is nonpositive. This gives Proposition 3's first estimate for every nonnegative step, with its actual estimator-energy term.



\noindent Lean declarations: \texttt{sgld\_one\_step}.

\section{Obtain the RMS recurrence and propagate finite transport}


When hM is at most two, drop the nonpositive score term. The exact component RMS bounds estimator energy by R squared times the initial squared radius about the minimizer. The radius estimate is twice the squared transport error plus 2p/m, yielding exactly the second claim's coefficients. Actual proposals preserve L2 and probability, so finite quadratic transport propagates along every iteration.



\noindent Lean declarations: \texttt{sgld\_one\_step\_rms}, \texttt{sgldLaws\_finite}.

\section{Prove the constant-step main theorem}


The full gradient has Lipschitz constant at most the exact R=n M\_rms, and strong convexity gives m at most R. The paper's step restriction therefore ensures the local stability condition. Apply the proved scalar recurrence to the actual SGLD marginal laws. The result has the reciprocal contraction and the precise persistent error A+2m+2 R squared/m. Substituting m=n mbar and dividing the full average by n produces the printed theorem in average-potential units.



\noindent Lean declarations: \texttt{sgld\_constant\_step}, \texttt{sgld\_constant\_step\_average}.

\section{Remove a finite-moment hypothesis from the main statement}


The extended endpoint applies the positive-affine transport extension. It has the same contraction, step restriction and average-potential constants for every probability initialization. For infinite initial distance its right side is infinite; otherwise all moments used by the real proof are derived from finite transport to the Gibbs target. The source-facing paper endpoint uses the literal coordinate average of the divided potential f/n. Gradient linearity rescales every coordinate Taylor quotient, and the proved homogeneity of their exact suprema identifies this constant with the full average divided by n, without a chosen replacement bound.



\noindent Lean declarations: \texttt{sgld\_constant\_step\_extended}, \texttt{sgld\_constant\_step\_paper}.

\end{document}
